\documentclass[11pt]{article}

% basic packages
\usepackage[margin=1in]{geometry}
\usepackage{amsmath, amssymb}
\usepackage{graphicx}
\usepackage{float}
\usepackage{booktabs}
\usepackage{hyperref}

% title info
\title{EE/CS 148B Homework 3 Writeup}
\author{Your Name}
\date{\today}

\begin{document}

\maketitle

\section*{1. Assignment Overview}

This writeup contains my responses, results, plots, and discussion for Homework 3.

\section*{2. Vision Transformer}

\subsection*{2.4 Design Questions}

\subsubsection*{CLS Token vs. Mean Pooling}

Write your 3--4 sentence response here.

\subsubsection*{Effect of Patch Size}

For a $224 \times 224$ image, the number of patches is

\[
N = \left(\frac{224}{P}\right)^2.
\]

\begin{center}
\begin{tabular}{c c c}
\toprule
Patch Size $P$ & Number of Patches $N$ & Time Mean $\pm$ Std \\
\midrule
8  & 784 & insert result here \\
16 & 196 & insert result here \\
32 & 49  & insert result here \\
\bottomrule
\end{tabular}
\end{center}

Write your 2--3 sentence discussion here.

\section*{3. CLIP-Style Contrastive Pretraining}

\subsection*{3.2 Symmetric InfoNCE Loss}

The loss is symmetric because it trains image-to-text matching and text-to-image matching at the same time.

\subsection*{3.3 Pretraining}

\subsubsection*{Training Loss Curve}

\begin{figure}[H]
    \centering
    \includegraphics[width=0.75\textwidth]{figures/training_loss.png}
    \caption{Training loss over epochs.}
\end{figure}

\subsubsection*{Zero-Shot Validation Accuracy Curve}

\begin{figure}[H]
    \centering
    \includegraphics[width=0.75\textwidth]{figures/val_accuracy.png}
    \caption{Zero-shot validation accuracy over epochs.}
\end{figure}

Write your 2--3 sentence discussion here.

\subsection*{3.4 Qualitative Analysis}

\begin{center}
\begin{tabular}{c c c}
\toprule
Example & True Label & Predicted Label \\
\midrule
1 & insert label & insert label \\
2 & insert label & insert label \\
3 & insert label & insert label \\
4 & insert label & insert label \\
5 & insert label & insert label \\
6 & insert label & insert label \\
7 & insert label & insert label \\
8 & insert label & insert label \\
9 & insert label & insert label \\
10 & insert label & insert label \\
\bottomrule
\end{tabular}
\end{center}

Write your 3--4 sentence discussion here.

\section*{4. LoRA Fine-Tuning}

\subsection*{4.1 LoRA-Wrapped Linear Layer}

Total parameters: insert value here.

Trainable parameters: insert value here.

Trainable ratio: insert value here.

\subsection*{4.2 Comparing Adaptation Strategies}

\begin{center}
\begin{tabular}{l c c c c}
\toprule
Method & Test Accuracy & Trainable Params & Peak Memory & Time \\
\midrule
Linear Probe & insert & insert & insert & insert \\
LoRA & insert & insert & insert & insert \\
Full Fine-Tuning & insert & insert & insert & insert \\
\bottomrule
\end{tabular}
\end{center}

Write your 4--5 sentence discussion here.

\subsection*{4.3 Rank Sweep}

\begin{figure}[H]
    \centering
    \includegraphics[width=0.75\textwidth]{figures/lora_rank_sweep.png}
    \caption{Test accuracy as a function of LoRA rank.}
\end{figure}

Write your 3--4 sentence discussion here.

\section*{5. Vision-Language Model}

\subsection*{5.3 Vision-Language Projector}

We need more than one linear layer because the projector must learn a richer mapping between the vision encoder space and the decoder embedding space. The extra nonlinearity gives the model more capacity while the encoder and decoder are frozen.

\subsection*{5.4 Token Injection Strategies}

\begin{center}
\begin{tabular}{l c c c c}
\toprule
Injection Strategy & Accuracy & Visual Tokens & Peak Memory & Time per Step \\
\midrule
CLS-only & insert & insert & insert & insert \\
All-patches & insert & insert & insert & insert \\
Interleaved & insert & insert & insert & insert \\
\bottomrule
\end{tabular}
\end{center}

Write your discussion here.

\subsection*{5.5 Attention Masking}

\subsubsection*{Mask Diagrams}

Insert your two mask diagrams here.

\begin{center}
\begin{tabular}{l c}
\toprule
Mask Type & Validation Accuracy \\
\midrule
Fully Causal & insert \\
Image Bidirectional & insert \\
\bottomrule
\end{tabular}
\end{center}

Write your 3--4 sentence discussion here.

\subsection*{5.6 Freezing Strategies}

\begin{center}
\begin{tabular}{l c c c}
\toprule
Configuration & Validation Accuracy & Trainable Params & Peak Memory \\
\midrule
A: Projector Only & insert & insert & insert \\
B: Projector + Decoder LoRA & insert & insert & insert \\
C: Projector + Full Decoder & insert & insert & insert \\
D: Full Fine-Tuning & insert & insert & insert \\
\bottomrule
\end{tabular}
\end{center}

Write your 5--6 sentence discussion here.

\subsection*{5.7 Qualitative Evaluation}

\begin{center}
\begin{tabular}{c p{4cm} c c}
\toprule
Example & Question & Ground Truth & Model Output \\
\midrule
1 & insert question & insert & insert \\
2 & insert question & insert & insert \\
3 & insert question & insert & insert \\
4 & insert question & insert & insert \\
5 & insert question & insert & insert \\
6 & insert question & insert & insert \\
7 & insert question & insert & insert \\
8 & insert question & insert & insert \\
9 & insert question & insert & insert \\
10 & insert question & insert & insert \\
\bottomrule
\end{tabular}
\end{center}

Write your 4--5 sentence discussion here.

\section*{6. Positional Encodings and RoPE}

\subsection*{6.1 1D RoPE}

Measured norm difference after applying RoPE: insert value here.

\subsection*{Learned PE vs. RoPE}

\begin{center}
\begin{tabular}{l c c}
\toprule
Positional Encoding & Train-Size Accuracy & Extrapolated-Size Accuracy \\
\midrule
Learned PE & insert & insert \\
1D RoPE & insert & insert \\
\bottomrule
\end{tabular}
\end{center}

Write your 3--4 sentence discussion here.

\subsection*{6.2 2D RoPE}

\begin{center}
\begin{tabular}{l c c}
\toprule
Positional Encoding & Train-Size Accuracy & Extrapolated-Size Accuracy \\
\midrule
1D RoPE & insert & insert \\
2D RoPE & insert & insert \\
\bottomrule
\end{tabular}
\end{center}

Write your 2--3 sentence discussion here.

\subsection*{6.3 Multimodal RoPE}

\subsubsection*{Naive 1D Position IDs}

Write paragraph 1 here.

\subsubsection*{First Text Token Position under M-RoPE}

Write paragraph 2 here.

\subsubsection*{Why M-RoPE Uses Three Chunks}

Write paragraph 3 here.

\subsection*{Bonus: M-RoPE Implementation}

\begin{center}
\begin{tabular}{l c c}
\toprule
Position Method & Overall Accuracy & Spatial-Question Accuracy \\
\midrule
Naive 1D & insert & insert \\
M-RoPE & insert & insert \\
\bottomrule
\end{tabular}
\end{center}

Write your 3--4 sentence discussion here.

\end{document}